%%vsgtu709

%

%

%%%%

%%%   Самарский государственный технический университет

%%%   Кафедра прикладной математики и информатики

%%%

%%%   Преамбула для статьи в журнал "Вестник Самарского государственного

%%%   технического университета. Серия: «Физико-математические науки»"

%%%   Ориентирована под MikTEx 2.4 for Win32

%%%

%%%   Хотя сам журнал собирается под GNU/Linux на дистрибутиве TeXLive



\documentclass[11pt,a4paper,twoside]{article}



\usepackage[cp1251]{inputenc}

\usepackage[T2A]{fontenc}

\usepackage[english,russian]{babel}

\usepackage{samgtubib}

\usepackage[dvips]{graphicx}

\usepackage[multiple]{footmisc}



\usepackage{samgtu}

\renewcommand{\VestnikYear}{2009} % Год издания Вестника

\renewcommand{\VestnikNumber}{2\,(19)} % Номер Вестника

\renewcommand{\VestnikMonth}{Ноябрь} % Месяц выхода Вестника

\renewcommand{\VestnikData}{01.11.09} % Дата подписания Вестника в печать

\renewcommand{\VestnikUPL}{26,64} % Усл. печ. л.

\renewcommand{\VestnikUIL}{25,73} % Уч.-изд. л.

\renewcommand{\VestnikRegNumber}{479/09} % Регистрационный номер

\renewcommand{\VestnikOrder}{994} % Номер заказа







%%

% \begin{document}



\renewcommand{\firstpage}{189}

\renewcommand{\lastpage}{198}





%%%%%%%%%%%%%%%%%%%%%%%%%%%%%%%%%%%%%%%%%%%%%%%%%%%%%%%%%%%%%%%%%%%%%%%%%%%%%%%%%%%%

%Информация о статье и об авторах на русском и английском языках



% Номер УДК

\udk{530.161+532.5+536.7}%Обработка текста. Подготовка текстов : Языки программирования. Компьютерные языки

\msc{76A02}% Коды MSC см. http://www.ams.org/msc/



\title{Термодинамика вязкой жидкости с~точки зрения наблюдателя}% Полный заголовок

{Термодинамика вязкой жидкости с~точки зрения наблюдателя}% Заголовок, который идёт в верхний колонтитул



\author{М.\,Ю.~Белевич}% Автор(ы) для заголовка, если авторов несколько укать \inst

{М.~Ю.~Б\,е\,л\,е\,в\,и\,ч}% Автор(ы) для оглавления и колонтитула через \,



\institute{\spbiok.}

% Если несколько, то так  \institute{ Организация 1 \and Организация 2  \and Организация 3}



\email{E-mail}{mbelevich@yahoo.com} %E-mail's авторов



\otherinf{\fullauthor{Михаил Юрьевич Белевич} \regalia{к.ф.-м.н., доц.}\position{ведущий  научный сотрудник} \dept{лаб. геофизических пограничных слоев} }







\keywords{неравновесная термодинамика, механика жидкости, законы сохранения, причинность}



\workabstract{Рассматривается вариант построения неравновесной термодинамики вязкой

жидкости, не требующий привлечения гипотезы локального термодинамического

равновесия. Предлагаемая теория является частью причинно обусловленной

модели вязкой теплопроводящей жидкости, включающей наблюдателя как

элемент описания движения среды. Ключевой момент "---  замена первого

начала термодинамики законом сохранения полной энергии "---  теоремой,

следующей из закона сохранения массы. Обсуждаются условия применимости

второго начала термодинамики и проблема диссипации кинетической энергии.

Основные выводы иллюстрируются примерами из численного анализа.}



\engtitle{Thermodynamics of the viscous fluid from an observer's viewpoint}



\engauthor{M.\,Yu.~Belevich}{B\,e\,l\,e\,v\,i\,c\,h~M.\,Yu.}





\enginstitute{\engspbiok.}% Место работы каждого из авторов на англ. языке

% Если несколько, то так  \enginstitute{ Организация 1 \and Организация 2  \and Организация 3}



\otherenginfm{\fullauthor{Michael Yu. Belevich} \regalia{Ph.\,D. (Phys. \& Math.)} \position{Leading Research Associate} \dept{Lab. of Geophysical Boundary Layers} } 





\engabstract{The development of the non-equilibrium thermodynamics of the viscous fluid not using the local thermodynamic equilibrium hypothesis is considered. The theory is based on the causal mechanics of the heat conducting continuum, which includes the 1st law of thermodynamics as a theorem. The conditions of applicability of the 2nd law of thermodynamics and the dissipation of the kinetic energy problem are discussed. Main conclusions are illustrated using the example from the numerical analysis. }



\engkeywords{non-equilibrium thermodynamics, fluid mechanics, conservation

laws, causality}

\date{01/X/2012} % Дата поступления статьи в редакцию.

\revisiondate{12/I/2013} % В окончательном виде



%%%%%%%%%%%%%%%%%%%%%%%%%%%%%%%%%%%%%%%%%%%%%%%%%%%%%%%%%%%%%%%%%%%%%%%%%%%%%%%%%%%%





\maketitle



%Каждый пункт статьи должен начинаться с команды Название пункта

%после названия точку ставить не нужно. Пункты автоматически нумеруются.

%Если пункт нумеровать не нужно, то следует использовать в команде

% \Section необязательный параметр [n]:

%   \Section[n]{Имя пункта}

% Введение и заключение обычно не нумеруются.





\Section[n]{Введение}

Стандартная модель вязкой жидкости включает в себя две теории. Первая "--- 

описывает движение жидкости под действием механических сил и основывается

на двух интегральных законах: сохранения массы и движения. Вторая

теория призвана описывать эволюцию характеристик жидкости под действием

термических факторов. Однако в классическом случае фундаментальные

законы, на основе которых можно описывать такую эволюцию, отсутствуют.

Имеется набор так называемых начал равновесной термодинамики, которые

распространяются на общий неравновесный случай с помощью гипотезы

локального термодинамического равновесия. Эта гипотеза предъявляет

ряд требований к изучаемым явлениям, и~для их преодоления были разработаны

специальные расширения классической термодинамики (см., например,

[\citen{MR93}--\citen{Jou99}]). 



Помимо указанных расширений классической термодинамики имеется \linebreak иной

путь построения общей термодинамической теории, применимой в неравновесном

случае. Он был разработан как часть причинно обусловленной модели

жидкости, предложенной в [\citen{Belevich03}--\citen{Belevich08}].

Эта модель допускает естественное рассмотрение термодинамических эффектов

без привлечения законов равновесной термодинамики и гипотезы локального

равновесия, в~то время как стандартная модель жидкости помимо постулатов

механики основана на постулате о сохранении полной энергии. Причинная

модель рассматривает сохранение энергии как теорему, следующую из

закона сохранения массы. По этой причине здесь не требуется ни постулирования

первого начала термодинамики, ни принятия гипотезы локального равновесия.



Целью данной работы является демонстрация возможности указанного способа

построения неравновесной термодинамики. Теория строится с точки зрения

наблюдателя, который является одним из элементов описания. Легко показать,

что стандартная гидромеханика является частным случаем этой более

общей причинно обусловленной теории. Настоящая работа имеет следующую структуру.

Раздел 1 содержит краткое описание причинной

модели вязкой жидкости, предложенной в \cite{Belevich03}. Баланс

энергии в причинной модели рассматривается в разделе~2,

а его сравнение с классическим случаем "---  в~разделе~3.

Раздел 4 посвящён обсуждению второго закона термодинамики.

Заключительные замечания можно найти в разделе 5.



\smallskip



\Section{Причинно обусловленная модель жидкости\label{sec:Causal-Fluid-Model}}



{\tt Основные понятия.} Гипотеза сплошности считается принятой, и любое тело $\mathcal{B}$

мыслится как сплошная среда. Каждая точка тела ассоциируется с~множеством

событий, образующих её \emph{мировую линию $\phi$} в \emph{пространстве

событий} $\mathcal{W}$. Совокупность мировых линий точек тела образует

его мировую трубку, рассматриваемую как 4-мерное многообразие $\mathcal{B}^{4}$

в пространстве $\mathcal{W}$.



Каждая мировая линия может быть гладко параметризована вещественным

параметром. Параметризация произвольна, и для сужения произвола могут

быть наложены дополнительные ограничения. Цель "--- заменить бесконечное

число независимых параметров единственным параметром "--- \emph{временем}

$t$. Для этого выделяется одна из мировых линий, называемая \emph{мировой

линией наблюдателя} $\phi_{0}$. Ее параметризация произвольна. Остальные

параметры синхронизируются с временем наблюдателя. Метод синхронизации

удобно интерпретировать в терминах скорости сигнала, используемого

для проведения наблюдений. Фазовая скорость сигнала будет обозначаться

через $c$. Наблюдатель, таким образом, есть совокупность параметров $(\phi_{0},\, t,\, c)$.



Наличие наблюдателя позволяет определить \emph{пространства одновременных

событий} $\mathcal{W}_{t}$. Сечение $\mathcal{B}_{t}$ мировой трубки

$\mathcal{B}^{4}$ пространством $\mathcal{W}_{t}$ называется \emph{конфигурацией}

тела $\mathcal{B}$ в момент времени $t$. Разные наблюдатели порождают

различные пространства одновременных событий и тем самым "--- разные

конфигурации тела для одного и того же момента времени. 



\smallskip



{\tt Интегральные законы сохранения.}

На теле $\mathcal{B}$ и на пространстве событий $\mathcal{W}$ можно

определить меры $\mathcal{M}(\mathcal{B})$ и $\mathcal{V}(\mathcal{W})$,

которые на сечении $\mathcal{B}_{t}$ индуцируют меры $m(\mathcal{B}_{t})$

и $V(\mathcal{B}_{t})$, обычно интерпретируемые как масса и объём

конфигурации тела. Постулируется закон сохранения 

\begin{equation}

d_{t}m(\mathcal{B}_{t})=0.\label{eq:1}

\end{equation}





Нетрудно показать (см. \cite{Belevich03,Belevich08}), что меру $m(\mathcal{B}_{t})$

можно рассматривать и как энергию%

\footnote{Чтобы различать обе интерпретации, используются разные обозначения

для одной и той же меры. Здесь физическая размерность энергии $cm$

совпадает с размерностью массы $m$. Для получения величины с обычной

размерностью энергии используется метрика.%

} $cm(\mathcal{B}_{t})$ конфигурации $\mathcal{B}_{t}$, что приводит

к двоякой интерпретации закона сохранения \eqref{eq:1}. Действительно,

если для любой конфигурации $\mathcal{B}_{t}$ выполняется равенство

$cm(\mathcal{B}_{t})=C\, m(\mathcal{B}_{t})$, $C=\textnormal{const}$,

тогда уравнение

\begin{equation}

d_{t}cm(\mathcal{B}_{t})=0\label{eq:2}

\end{equation}

необходимо следует из уравнения \eqref{eq:1}. Реализация такой возможности

обсуждается ниже.



\smallskip



{\tt Дифференциальные законы сохранения.}

При весьма общих допущениях существует единственная функция $\rho$

такая, что масса и объём конфигурации оказываются связаны соотношением

$$m(\mathcal{B}_{t})=\int_{V(\mathcal{B}_{t})}\rho dV.$$ Аналогичное

соотношение связывает и энергию с объёмом: $$cm(\mathcal{B}_{t})=\int_{V(\mathcal{B}_{t})}\kappa dV.$$

Величины $\rho$ и $\kappa$ называются плотностями массы и энергии

в точке, и при этом $\kappa=C\rho$. 



Уравнения \eqref{eq:1} и \eqref{eq:2} порождают соответствующие дифференциальные законы сохранения

\begin{equation}

\textnormal{div}\left(\rho\vec{\tau}\right) =  0,\label{eq:3}

\end{equation}

$$\textnormal{div}\left(\kappa\vec{\tau}\right)=\rho d_{t}\frac{\kappa}{\rho} = 0,$$

где $\vec{\tau}=d_{t}$ "---  вектор, касательный к мировой линии.

Поскольку по определению $\kappa=\frac{1}{2}|\vec{\tau}|^{2}\rho$

и $\kappa=C\rho$, имеем $C=\frac{1}{2}|\vec{\tau}|^{2}$. За счёт

выбора подходящего метрического тензора $\mathsf{g}$ величину $|\vec{\tau}|^{2}=\mathsf{g}(\vec{\tau},\vec{\tau})$

можно положить равной константе, например, $|\vec{\tau}|=1$. 

В~этом случае $\kappa=\frac{1}{2}\rho$. 

Относительно ортогонального базиса метрический тензор диагонален и может быть выбран пропорциональным единичному тензору $\mathsf{I}$ с~коэффициентом пропорциональности

$g_{0}$: $\mathsf{g}=g_{0}\mathsf{I}$. При этом $g_{0}^{-1}=\mathsf{I}(\vec{\tau},\vec{\tau})$.



\smallskip



{\tt Осреднение.}

Поскольку гипотеза сплошности принята, движение далее следует разделить по масштабам и осреднить, т.~е. перейти к термодинамическому случаю.

При этом каждая мировая линия представляется в виде суммы сглаженной кривой и пульсации. Соответственно, касательное векторное поле $\vec{\tau}=\vec{v}+\vec{\tau}'$

разделяется на осреднённое векторное поле $\vec{v}$, касающееся сглаженных

мировых линий, и векторное поле пульсаций $\vec{\tau}'$. Обсуждение

процедуры осреднения см. в \cite{Belevich04,Belevich09}.



В результате осреднения (далее обозначаемого надчеркиванием) закон

сохранения массы \eqref{eq:3} приводит к уравнению диффузии плотности

массы, которое следует рассматривать как обобщение уравнения неразрывности

(детали см. в \cite{Belevich09}). Однако как правило,  рассматривают

более грубое приближение, сохраняющие вид дифференциального закона

сохранения массы. Здесь осредненное уравнение \eqref{eq:3} также

будем записывать в виде 4-мерного уравнения неразрывности

\[

\overline{\textnormal{div}\left(\rho\vec{\tau}\right)}=\textnormal{div}(\rho\vec{v})=0.

\]





В свою очередь, энергия (теперь называемая \emph{полной энергией})

записывается как сумма \emph{кинетической} $K(\mathcal{B}_{t})$ и

\emph{внутренней} $E(\mathcal{B}_{t})$ энергий конфигурации тела

$\overline{cm(\mathcal{B}_{t})}=K(\mathcal{B}_{t})+E(\mathcal{B}_{t}),$

$K,\, E>0$ с плотностями 

\begin{equation}

\bar{\kappa}\equiv\frac{1}{2}\rho\left|\vec{v}\right|^{2},\qquad\varepsilon\equiv\frac{1}{2}\rho\bigl(\overline{|\vec{\tau}|^{2}}-|\vec{v}|^{2}\bigr)=\frac{1}{2}\rho\left(1-|\vec{v}|^{2}\right).

\label{eq:5}

\end{equation}

В силу определения \eqref{eq:5} верно равенство $$\frac{\bar{\kappa}+\varepsilon}{\rho}=\frac{1}{2}.$$

Метрический коэффициент $g_{0}$ в этом случае имеет вид 

$$g_{0}=\frac{1-2\epsilon}{\mathsf{I}(\vec{\tau},\vec{\tau})},\quad \epsilon=\frac{\varepsilon}{\rho}.$$ 

Осредняя скорость изменения энергии, получим

\begin{eqnarray*}

\overline{d_{t}cm} & = & d_{t}\underbrace{\int_{V}\left(\bar{\kappa}+\varepsilon\right)dV}_{(K+E)}+\underbrace{\int_{\partial V}(\mathsf{T}'\vec{v})_{n}dS}_{=0}+\underbrace{\int_{\partial V}(\rho\overline{\epsilon'\vec{\tau}'})_{n}dS}_{=Q}=0.

\end{eqnarray*}

Первое слагаемое описывает скорость изменения полной энергии. Интеграл

$$\int_{\partial V}(\mathsf{T}'\vec{v})_{n}dS$$ равен нулю, поскольку

нормальная компонента $\left(\mathsf{T}'\vec{v}\right)_{n}$ потока

через границу $\partial V$ сечения мировой трубки $\mathcal{B}_{t}$

отсутствует. Последнее слагаемое $$\int_{\partial V}(\rho\overline{\epsilon'\vec{\tau}'})_{n}dS$$

описывает внешний поток $Q$ внутренней энергии через границу. Здесь

величина $$\epsilon'\equiv\frac{1}{2}|\vec{\tau}'|^{2}=\frac{\varepsilon'}{\rho}$$

интерпретируется как удельная плотность пульсаций внутренней энергии,

порождаемых пульсациями скорости $\vec{\tau}'$. Величина $\rho\overline{\vec{\tau}'\otimes\vec{\tau}'}=\mathsf{T'}$ "--- 

4-мерный \emph{тензор вязких напряжений}, а $\mathsf{T}'\vec{v}$

означает свертку $\mathsf{T'}$ с $\vec{v}$. В отсутствие $Q$ имеет

место закон сохранения $$d_{t}(K+E)=0,$$ которому соответствует дифференциальное

уравнение баланса плотности полной энергии

\begin{equation}

0=\rho\underbrace{d_{t}\frac{\bar{\kappa}+\varepsilon}{\rho}}_{=0}+\textnormal{div}(\mathsf{T}'\vec{v}+\rho\overline{\epsilon'\vec{\tau}'}).\label{eq:7-2}

\end{equation}





\smallskip



\Section{Баланс энергии\label{sec:Energy-balance}}







{\tt Уравнения баланса.}

Из уравнения \eqref{eq:7-2} можно получить уравнения баланса плотности

обоих видов энергии. Продифференцируем слагаемое $$\textnormal{div}(\mathsf{T}'\vec{v})=\mathsf{g}(\vec{v},\textnormal{div}\mathsf{T}')+\mathsf{T}'\colon\nabla\vec{v},$$

а затем сгруппируем члены в уравнении \eqref{eq:7-2} следующим образом:

\begin{equation}

\underbrace{\rho d_{t}\frac{\bar{\kappa}}{\rho}+\mathsf{g}(\vec{v},\textnormal{div}\mathsf{T}')}_{=d_{t}\pi}+\underbrace{\rho d_{t}\epsilon+\mathsf{T}'\colon\nabla\vec{v}+\textnormal{div}\rho\overline{\epsilon'\vec{\tau}'}}_{=-d_{t}\pi}=0.\label{eq:te1}

\end{equation}

Здесь $(\cdot):(\cdot)$ "---  скалярное произведение тензоров. Поскольку

обе группы членов должны различаться лишь знаком, обозначим их через

$\pm d_{t}\pi$ и запишем два дифференциальных уравнения баланса энергии

\begin{eqnarray}

\rho d_{t}\frac{\bar{\kappa}}{\rho} & = & d_{t}\pi-\mathsf{g}(\vec{v},\textnormal{div}\mathsf{T}'),\label{eq:ke1}\\

\rho d_{t}\epsilon & = & -d_{t}\pi-\mathsf{T}'\colon\nabla\vec{v}-\textnormal{div}\rho\overline{\epsilon'\vec{\tau}'},\label{eq:ie1}

\end{eqnarray}

соответственно кинетической и внутренней. Величина\textit{\emph{ $\pi$

имеет смысл скалярной кривизны пространства событий $\mathcal{W}$

(детали см. в \cite{Belevich05}).}} Слагаемое $\mathsf{g}(\vec{v},\textnormal{div}\mathsf{T}')$

в уравнении \eqref{eq:ke1} описывает перераспределение и диссипацию

плотности кинетической энергии, тогда как слагаемые $\mathsf{T}'\colon\nabla\vec{v}$

и $\textnormal{div}\rho\overline{\epsilon'\vec{\tau}'}$ в уравнении

\eqref{eq:ie1} "---  вязкий приток внутренней энергии и ее перенос

теплопроводностью.



\smallskip

{\tt Баланс импульса.}

Слагаемые в уравнении \eqref{eq:te1} были сгруппированы так, чтобы

первая группа соответствовала уравнению движения вязкой жидкости.

Действительно, скорость изменения удельной плотности кинетической

энергии можно записать в виде

\begin{equation}

\rho d_{t}\frac{\bar{\kappa}}{\rho}=\mathsf{g}(\vec{v},\textnormal{div}\mathsf{M}),\label{eq:ke1a}

\end{equation}

где тензор $\mathsf{M}\equiv\rho\vec{v}\otimes\vec{v}$ называется

4-\textit{тензором плотности потока импульса}. Постулируемое уравнение

\begin{equation}

\textnormal{div}\mathsf{M}=\textnormal{div}\mathsf{T},\label{e13}

\end{equation}

где $\mathsf{T}$ "---  \emph{4-тензор напряжений}\textit{\emph{, известно

как уравнение баланса плотности импульса или уравнение движения}}

(в данном случае 4-мерное). В соответствии с \eqref{e13} уравнение

\eqref{eq:ke1a} можно переписать так:

\begin{equation}

\rho d_{t}\frac{\bar{\kappa}}{\rho}=\mathsf{g}(\vec{v},\textnormal{div}\mathsf{T}).\label{eq:ke2}

\end{equation}





Комбинируя \eqref{eq:7-2}, \eqref{e13} и \eqref{eq:ke2}, получим

уравнение баланса плотности внутренней энергии (ср. с уравнением \eqref{eq:ie1})

\begin{equation}

\rho d_{t}\epsilon=-\mathsf{g}(\vec{v},\textnormal{div}\mathsf{T})-\textnormal{div}\rho\overline{\epsilon'\vec{\tau}'}.\label{eq:ie2}

\end{equation}



\smallskip



{\tt Тензор напряжений.}

\textit{\emph{Сравнивая оба варианта записи уравнения баланса плотности

кинетической энергии}} \eqref{eq:ke1} и \eqref{eq:ke2}\textit{\emph{,

получим

\begin{equation}

\mathsf{T}=\pi\mathsf{g}^{-1}-\mathsf{T}'.\label{eq:stress}

\end{equation}

Тензор вязких напряжений определяется как}}\emph{ }\textit{\emph{$\mathsf{T}'=2\rho\chi\mathsf{D}$.

Здесь $\mathsf{D}$ "---  4-мерный }}\textit{тензор скоростей деформации}\textit{\emph{.

Величина $\mu=-\rho\chi g_{0}^{-1}$ называется }}\textit{коэффициентом

динамической вязкости}\textit{\emph{. Первое слагаемое в правой части

\eqref{eq:stress} обычно записывается в виде $\pi\mathsf{g}^{-1}=\pi g_{0}^{-1}\mathsf{I}=-p\mathsf{I},$

а множитель $p\equiv-\pi g_{0}^{-1}$ называется }}\textit{давлением}\textit{\emph{.

Таким образом, вид 4-мерного тензора напряжений $\mathsf{T}$, записанного

в терминах давления и тензора скоростей деформации, согласуется с его

3-мерным аналогом $\mathsf{T}=-p\mathsf{I}+2\mu\mathsf{D}$. Если

пренебречь тензором вязких напряжений, получим модель идеальной жидкости,

в противном случае "---  модель вязкой жидкости.}}



\smallskip



{\tt Уравнение баланса внутренней энергии.}

Слагаемое $\textnormal{div}\rho\overline{\epsilon'\vec{\tau}'}$ в

правой части уравнения \eqref{eq:ie2} описывает перенос плотности

внутренней энергии теплопроводностью. Соответствующая корреляция обычно

записывается в терминах градиента удельной плотности внутренней энергии

$\overline{\epsilon'\vec{\tau}'}=\theta\mathsf{g}^{-1}(\nabla\epsilon)$.

То же в случае ортогонального базиса выглядит так: $\overline{\epsilon'\vec{\tau}'}=-k\nabla\epsilon$,

где $k=-g_{0}^{-1}\theta$. Таким образом, рассматриваемое слагаемое

теперь имеет вид $\textnormal{div}\rho\overline{\epsilon'\vec{\tau}'}= $ \linebreak $\textnormal{div}\left(\rho\theta\mathsf{g}^{-1}(\nabla\epsilon)\right)$

или $\textnormal{div}\rho\overline{\epsilon'\vec{\tau}'}=-\textnormal{div}\left(\rho k\nabla\epsilon\right)$,

если базис ортогонален. Множитель $k$ называется \emph{коэффициентом

температуропроводности}. 



В простейшем случае тензор напряжений $\mathsf{\mathsf{T}=\pi\mathsf{g}^{-1}}$.

Уравнение баланса внутренней энергии \eqref{eq:ie2} в случае идеальной

жидкости имеет вид $\rho d_{t}\epsilon=$ \linebreak $= -d_{t}\pi-\textnormal{div}\rho\overline{\epsilon'\vec{\tau}'}$.

Полагая \textsf{$\mathsf{g}=g_{0}\mathsf{I}$} и вводя обозначение

$e=g_{0}^{-1}\epsilon$ для удельной плотности внутренней энергии,

получим уравнение 

\begin{equation}

\rho d_{t}e=d_{t}p+\textnormal{div}\left(\rho k\nabla e\right).\label{Intbalance}

\end{equation}





В более общем случае $\mathsf{T}=\pi\mathsf{g}^{-1}+2\rho\chi\mathsf{D}$,

и\textsf{ }если \textsf{$\mathsf{g}=g_{0}\mathsf{I}$}, получим аналогичное

уравнение 

\begin{equation}

\rho d_{t}e=d_{t}p-\textnormal{div}\left(2\mu\mathsf{D}\vec{v}\right)+2\mu\mathsf{D}:\mathsf{D}+\textnormal{div}\left(\rho k\nabla e\right).\label{eq:4D_Int}

\end{equation}



\smallskip



\Section{Сравнение со стандартной моделью жидкости\label{sub:Comparison}}

В стандартной гидромеханике закон сохранения полной энергии "---  постулат.

Изначально он формулируется в рамках равновесной термодинамики и

его использование в рамках нестационарной модели среды возможно, если

выполняется гипотеза \emph{локального термодинамического равновесия}.

Понятие равновесия требует введения функции состояния $T$ "---  \emph{температуры.}

Утверждается, что две термодинамические системы находятся в равновесии,

если их температуры равны. 



В рассмотренной причинно обусловленной модели потребность в классической

термодинамике и гипотезе равновесия отсутствует. Это связано с тем,

что закон сохранения полной энергии в причинной модели является теоремой,

следующей из сохранения массы. При этом причинно обусловленное уравнение

баланса внутренней энергии отличается от классического. Сравнение

стандартного уравнения с классическим пределом ($c\to\infty$) соответствующего

причинного уравнения показывает, в чем состоит принятие гипотезы локального

равновесия.



\smallskip



{\tt Определение температуры.}

Дадим определение температуры $T$, применимое в общем неравновесном

случае. Разумно выбрать его таким, чтобы справедливыми оставались

известные результаты. Определим температуру $T$ с~помощью соотношения

$$de=c_{p}dT,$$ где множитель $c_{p}$ называется \emph{удельной теплоёмкостью}

среды \emph{при постоянном давлении}. Подставив это выражение в уравнение

\eqref{eq:4D_Int}, получим 4-мерное \emph{уравнение теплопроводности}

движущейся среды

\[

\rho c_{p}d_{t}T=d_{t}p-\textnormal{div}\left(2\mu\mathsf{D}\vec{v}\right)+2\mu\mathsf{D}:\mathsf{D}+\textnormal{div}\left(\lambda\nabla T\right).

\]

Оно является гиперболическим и было предложено в \cite{Belevich04}.

Здесь $\lambda\equiv\rho kc_{p}$ "---  \emph{коэффициент теплопроводности}. 



\smallskip



{\tt Локальное равновесие в жидкости.}

Стандартное уравнение баланса удельной внутренней энергии $e_{3}$

идеальной жидкости без учёта внешних источников записывается в виде

$$\rho d_{t}e_{3}=\frac{p_{3}}{\rho}d_{t}\rho.$$ Соответствующее причинное

уравнение \eqref{Intbalance} в классическом пределе ($c\to\infty$)

таково:

\[

\rho d_{t}e_{3}=d_{t}p_{3}\quad\Rightarrow\quad\rho d_{t}e_{3}=\frac{p_{3}}{\rho}d_{t}\rho+\rho d_{t}\frac{p_{3}}{\rho}.

\]

Здесь введены обозначения $e_{3}=\lim_{c\to\infty}e$ и $p_{3}=\lim_{c\to\infty}p$.

Легко видеть, что оба уравнения совпадают, если $$d_{t}\frac{p_{3}}{\rho}=0.$$

Это равенство следует рассматривать как математическое выражение

гипотезы локального равновесия в идеальной жидкости. Действительно,

для уравнения состояния $p_{3}=\rho RT$, где $R$ "---  универсальная

газовая постоянная, отношение $\frac{p_{3}}{\rho}$ равно $RT$ и

из $d_{t}\frac{p_{3}}{\rho}=0$ следует $T=\textnormal{const}.$



В случае вязкой жидкости аналогичными рассуждениями получим равенство

$$\rho d_{t}\frac{p_{3}}{\rho}-\textnormal{div}\left(2\mu\mathsf{D}_{3}\vec{v}\right)=0,$$

означающее выполнение условия локального равновесия в вязкой жидкости.



\smallskip



{\tt Диссипация кинетической энергии.}

В стандартной теории слагаемое $$\textnormal{diss}_{3}\equiv2\mu\mathsf{D}_{3}:\mathsf{D}_{3},$$

описывающее \emph{диссипацию кинетической энергии}, знакопостоянно.

В причинной модели жидкости знак аналогичного слагаемого $$\textrm{diss}=2\mu\mathsf{D}:\mathsf{D}$$

определяется соотношением, связывающим 3-мерную скорость $\vec{v}_{3}$ и скорость сигнала. 

Можно показать (детали см. в \cite{Belevich05}), что если

$|\vec{v}_{3}|^{2}<c^{2}$, то плотность диссипации будет положительна.



Скорость изменения кинетической энергии сечения мировой трубки имеет

вид 

\begin{eqnarray}

d_{t}(g_{0}^{-1}K) & = & -\int_{V(\mathcal{B}_{t})}d_{t}pdV-\textnormal{Diss},\label{eq:dissK}

\end{eqnarray}

где $$\textnormal{Diss}=\int_{V(\mathcal{B}_{t})}\textnormal{diss}\, dV$$ "--- 

потери кинетической энергии за счёт диссипации. Величина $g_{0}^{-1}K$ "--- 

кинетическая энергия тела, имеющая обычную физическую размерность.

Первое слагаемое в правой части уравнения \eqref{eq:dissK} описывает

изменение $K$ за счёт сжимаемости среды, а второе слагаемое "--- 

<<диссипацию>> кинетической энергии сечения. Поскольку часть слагаемых

в <<диссипации>> отрицательны, в некоторых специальных случаях <<диссипация>>

может восприниматься как <<антидиссипация>>. Интерпретация подобных

случаев дана ниже.



\smallskip



\Section{Второй закон термодинамики в общем случае\label{sec:The-2nd-law}}



{\tt Определение энтропии.}

Дадим для общего неравновесного случая определение плотности энтропии

$s$, которое не требует допущения локального термодинамического равновесия.

По аналогии с уравнением Гиббса (см., например, \cite{Gy74}) определим

плотность энтропии $s$ с помощью соотношения $$T\rho d_{t}\frac{s}{\rho}\equiv\rho d_{t}e-d_{t}p.$$

В случае локального равновесия $d_{t}\frac{p_{3}}{\rho}=0$ оно совпадает 

со стандартным определением и для идеальной жидкости равно нулю. Для

вязкой жидкости уравнение баланса плотности внутренней энергии даёт

уравнение баланса энтропии

\begin{eqnarray*}

T\rho d_{t}\frac{s}{\rho} & = & -\textnormal{div}\left(2\mu\mathsf{D}\left(\mathsf{g}(\vec{v})\right)\right)+\textrm{diss}.

\end{eqnarray*}



\smallskip



{\tt Второй закон термодинамики.}

Легко показать, что знак скорости изменения энтропии $d_{t}S$ определяется

знаком скалярного произведения $\mathsf{D}:\mathsf{D}$. В~стандартной

гидромеханике, как и в причинной теории при $c\rightarrow\infty$, выполняется

второй закон термодинамики и имеет место неравенство Клаузиуса $d_{t}S>0$.

В общем же причинном случае имеются две возможности: кроме указанной

стандартной возможности имеется и нестандартная "---  второй закон не

выполняется. 



Для того чтобы выяснить смысл этих двух возможностей, рассмотрим

частный случай $\vec{v}=\left(ic,u,0,0\right)$, $c=\textnormal{const}$.

Можно показать, что необходимым условием выполнения второго закона термодинамики

здесь является неравенство $$c^{2}>\frac{1}{2}\Bigl(\frac{\partial_{t}u}{\partial_{x}u}\Bigr)^{2}.$$

Если это неравенство выполняется, реализуется стандартная возможность:

наблюдается выполнение второго закона термодинамики, уменьшение кинетической

энергии за счёт диссипации и возрастание энтропии в замкнутой системе.

В противном случае сигнал оказывается слишком медленным и наблюдения

будут рисовать ошибочную картину явления, которая заключается в невыполнении

второго закона термодинамики и в поведении диссипативного слагаемого:

оно ведет себя как <<антидиссипация>> и приводит к росту кинетической

энергии. Всё это говорит о том, что используемый сигнал не подходит

для проведения наблюдений и должен быть заменён на более быстрый.



\pagebreak



\Section{Заключительные замечания\label{sec:Concluding-remarks}}



{\tt Общая схема рассуждений.}

Ключевой момент данного исследования состоит в замене постулата (закона

сохранения энергии) теоремой с тем же названием. В этом случае, кроме

выбора метрики, никаких дополнительных предположений делать не требуется.

Для того чтобы сформулировать и доказать указанную теорему, необходимо

рассмотреть 4-мерное пространство событий $\mathcal{W}$ и определить

в нём метрический тензор $\mathsf{g}$ исходя из условия \mbox{$|\vec{\tau}|=1$.}

При этом плотность энергии оказывается пропорциональной плотности

массы тела. Такая метрика может быть построена с помощью дополнительной

степени свободы $c$ "---  скорости сигнала, используемого для наблюдений.

Скорость сигнала ограничивает допустимые скорости наблюдаемых объектов.

Можно показать (см. \cite{Belevich03}), что наблюдение даёт правильное

значение скорости движущегося объекта тогда и только тогда, когда

скорость сигнала превосходит скорость объекта.



\smallskip



{\tt Второй закон термодинамики и численная неустойчивость.}

Вывод о том, что при недостаточно высокой скорости сигнала наблюдается

искаженная картина явления, непосредственно проверить не удается,

поскольку отсутствуют соответствующие эксперименты. Однако косвенная

проверка возможна. Речь идет о численном (здесь "---  конечно"=разностном)

решении уравнений модели. При численном рассмотрении соответствующих

задач сплошная среда заменяется набором узлов пространственно"=временной

сетки. Скорости используемых сигналов могут меняться в широком диапазоне от

конечных, равных отношению шагов сетки (явные модели), до бесконечных

(неявные модели). Анализ устойчивости таких моделей (см., например,

\cite{Roach}) показывает, что если скорость сигнала выбрана слишком

малой (меньшей, чем возможная скорость движения среды), вычисления

(здесь они играют роль наблюдений) дадут неверную картину движения

жидкости. Решение уравнений модели с помощью явных численных алгоритмов

подобно реальным наблюдениям движения жидкости, тогда как неявные

алгоритмы с их бесконечной скоростью сигнала соответствуют стандартной

механике жидкости.







%% Благодарности, гранты и прочее



\clear



\begin{thebibliography}{10}



\Bibitem{MR93}

\by I.~M\"uller, T.~Ruggeri

\book  Extended thermodynamics  

\serial Springer Tracts in Natural Philosophy

\vol 37

\publ Springer-Verlag 

\publaddr New York

\yr 1993

\totalpages xii+230





\Bibitem{Jou88}

\by D.~Jou, J.~Casas-V\'azquez, G.~Lebon

\paper Extended irreversible thermodynamics 

\jour  Rep. Prog. Phys. 

\vol  51

\issue 8

\pages 1105--1179 

\yr 1988







\Bibitem{Jou99}

\by D.~Jou, J.~Casas-V\'azquez, G.~Lebon

\paper  Extended irreversible thermodynamics revisited (1988--98) 

\jour  Rep. Prog. Phys. 

\vol 62 

\issue 7

\pages 1035--1142 

\yr 1999





\Bibitem{Belevich03}

\by M.~Belevich

\paper Causal description of non-relativistic dissipative fluid motion

\jour  Acta Mech.

\issue 1--2

\vol  161 

\pages 65--80

\yr 2003





\Bibitem{Belevich05}

\by M.~Belevich.

\paper Relationship between standard and causal fluid models

\jour  Acta Mech.

\issue 1--4

\vol  180

\pages 83--106

\yr 2005



\Bibitem{Belevich08}

\by M.~Belevich

\paper Non-relativistic abstract continuum mechanics and its possible physical interpretations

\jour  J.~Phys. A, Math. Theor. 

\vol 41

\issue 4

\yr 2008 

\papernumber 045401

\totalpages 19 





\Bibitem{Belevich04}

\by M.~Belevich

\paper Causal description of heat and mass transfer

\jour J.~Phys.~A, Math. Gen.

\vol  37

\issue 8

\pages 3053--3069 

\yr 2004



\Bibitem{Belevich09}

\by M.~Belevich 

\paper On the continuity equation

\jour   J.~Phys.~A, Math. Theor.

\vol 42

\issue 37

\yr 2009 

\papernumber 375502 

\totalpages 13



\Bibitem{Gy74} 

\by I.~Gyarmati 

\book  Non-equilibrium Thermodynamics. Field Theory and Variational Principles 

\publ Springer 

\publaddr New York 

\yr 1970



\Bibitem{Roach} 

\by P.~J.~Roach

\book  Computational Fluid Dynamics 

\publ Hermosa Publishers 

\publaddr Albuquerque, N.M. 

\yr 1976

\totalpages vii+446



\end{thebibliography}





\noindent

\hskip6.5cm \thedate \par\noindent

\hskip6.5cm\therevdate

%







%\chead[\fancyplain{}{\scriptsize \sl B\,e\,l\,e\,v\,i\,c\,h~M.\,Yu.}]{\fancyplain{}{\scriptsize \sl  Thermodynamics

%of the viscous fluid \dots}}





\makeengtitlem







\newpage

%

% \tableofengcontents

%

%\end{document}

